% 原书物理页=印刷页10--11；1.3 全节.
\section{有限 Abel 群}\label{sec:fg-1-3}
由循环群组成的（直）积，或更一般地由 Abel 群组成的（直）积，显然仍是 Abel 群. 其逆亦真：任何有限 Abel 群都同构于若干循环群的积. 更精确地说：
\begin{theorem}\label{thm:fg-1-3-1}
设 $G$ 为一个不等于 $\{0\}$ 的有限 Abel 群. 则存在唯一的一族整数 $a_1,\dots,a_s\geqslant2$，满足如下性质：
\begin{enumerate}[label=\roman*)]
\item 对每个 $i=1,\dots,s-1$，$a_i$ 整除 $a_{i+1}$；
\item $G$ 同构于积 $(\mathbb{Z}/a_1\mathbb{Z})\times\cdots\times(\mathbb{Z}/a_s\mathbb{Z})$.
\end{enumerate}
这些整数 $a_i$ 称为 $G$ 的\newterm{fg-invariant-factors}{不变因子}{invariant factors}.
\end{theorem}
\begin{proof}
下面的直接证明依赖于如下引理，其中有限群 $G$ 的指数指其各元素的阶的最小公倍数；等价地，即满足 $\forall x\in G,\ nx=0$ 的那些整数 $n$ 的最大公约数（或者又可说成：满足此条件的最小整数 $n\geqslant1$）.
\begin{lemma}\label{lemma:fg-1-3-2}
设 $G$ 为有限 Abel 群，$r$ 为 $G$ 的指数. 则 $G$ 含有阶为 $r$ 的元素.
\end{lemma}
\begin{proof}
用归纳法，只需证明：若 $x$ 与 $y$ 是 $G$ 中阶分别为 $n$ 与 $m$ 的两个元素，则 $G$ 中存在阶为 $r=\operatorname{lcm}(n,m)$ 的元素. 将 $n$ 与 $m$ 分解成素数幂之积，可见存在 $n$ 的因子 $n'$ 与 $m$ 的因子 $m'$，它们互素，且使得 $r=n'm'$. 于是 $x'=\frac{n}{n'}x$ 的阶为 $n'$，$y'=\frac{m}{m'}y$ 的阶为 $m'$，而 $x'y'$ 的阶为 $\operatorname{lcm}(n',m')=r$.
\end{proof}
\emph{各个 $a_i$ 的存在性（对 $|G|$ 作归纳）.}\ 设 $a_s$ 为 $G$ 的指数，$x_s$ 为 $G$ 中阶为 $a_s$ 的一个元素，并设 $G'$ 为商群 $G/\langle x_s\rangle$. 根据归纳假设，存在依次整除的整数 $a_1,\dots,a_{s-1}$，以及 $G'$ 中阶分别为 $a_1,\dots,a_{s-1}$ 的元素 $x'_1,\dots,x'_{s-1}$，使得 $G'=\langle x'_1\rangle\times\cdots\times\langle x'_{s-1}\rangle$；注意 $a_{s-1}$ 是 $G'$ 的指数，因而整除 $G$ 的指数 $a_s$. 我们来证明：对每个 $i=1,\dots,s-1$，存在 $G$ 中阶为 $a_i$ 的元素 $x_i$，使其在 $G'$ 中的像为 $x'_i$. 设 $y_i$ 是 $x'_i$ 在 $G$ 中的一个代表元. 则 $a_iy_i\in\langle x_s\rangle$，从而存在整数 $k_i$ 使得 $a_iy_i=k_ix_s$. 我断言 $a_i$ 整除 $k_i$，且 $x_i:=y_i-\frac{k_i}{a_i}x_s$ 满足要求. 事实上，$0=a_sy_i=\frac{a_s}{a_i}a_iy_i=\bigl(\frac{a_s}{a_i}k_i\bigr)x_s$，故 $a_s$ 整除 $\frac{a_s}{a_i}k_i$，且 $a_i\mid k_i$. 于是 $a_i\bigl(y_i-\frac{k_i}{a_i}x_s\bigr)=0$，故 $x_i$ 的阶整除 $a_i$，且由于 $x_i$ 在 $G'$ 中的像 $x'_i$ 已具有阶 $a_i$，故二者相等.

在这些条件下，$G$ 的子群 $H=\langle x_1\rangle\times\cdots\times\langle x_{s-1}\rangle$ 与 $K=\langle x_s\rangle$ 仅在 $0$ 处相交（$n_1x_1+\cdots+n_{s-1}x_{s-1}=n_sx_s\Rightarrow n_1x'_1+\cdots+n_{s-1}x'_{s-1}=0\Rightarrow a_1\mid n_1,\ \dots,\ a_{s-1}\mid n_{s-1}\Rightarrow n_1x_1=\cdots=n_{s-1}x_{s-1}=0$），且 $|G|=|G'|\cdot|K|=|H|\cdot|K|$. 于是 Abel 群 $G$ 是 $H$ 与 $K$ 的直积，从而得到所求的分解.
\emph{各个 $a_i$ 的唯一性.}\ 最大的那个不变因子就是 $G$ 的指数，故其唯一性是清楚的，但我们将从 $a_1$（以及 $s$）的唯一性入手. 对每个整数 $d$，令 $\Phi(d)=\operatorname{card}\{x\in G,\ dx=0\})=\prod_{i=1,\dots,s}\operatorname{gcd}(d,a_i)\leqslant d^s$\footnote{译注：原书在 $\operatorname{card}\{x\in G,\ dx=0\}$ 之后多印了一个右括号，此处照录.}，则 $s$ 是使得对一切 $d\geqslant1$ 都有 $\Phi(d)\leqslant d^k$ 的最小整数 $k$，因而只依赖于 $G$. 此外，$a_1$ 是使 $\Phi(d)\geqslant d^s$ 成立的最大整数 $d$，因而也只依赖于 $G$.

最后对 $|G|$ 作归纳以得出结论：设 $G[a_1]$ 为 $G$ 中由阶整除 $a_1$ 的元素组成的子群. 对每个 $i\geqslant1$，$(\mathbb{Z}/a_i\mathbb{Z})[a_1]=\frac{a_i}{a_1}\mathbb{Z}/a_i\mathbb{Z}$，因此 $G/G[a_1]$ 以 $\bigl(\frac{a_2}{a_1},\dots,\frac{a_s}{a_1}\bigr)$ 为其不变因子. 归纳假设推出它们只依赖于 $G$. 于是 $G$ 的不变因子 $(a_1,\dots,a_s)$ 亦然.
\end{proof}
